Copy mechanisms are fundamental to sequence-to-sequence models for summarization, dialogue, and code generation, yet the calibration of their copy-or-generate probability estimates has never been measured. We ask whether a small Transformer trained exclusively for copy decisions produces better-calibrated probabilities than a joint pointer-generator that learns generation and copying simultaneously. Using a synthetic copy task with known ground-truth labels, we compare a joint pointer-generator (824K parameters) against a specialized copy-decision model (472K parameters) across five random seeds. Both models achieve remarkably low expected calibration error (ECE $< 0.01$), with no statistically significant difference ($p = 0.27$). However, the specialized model dramatically outperforms the joint model in copy decision accuracy (F1: 0.985 vs.\ 0.969, a 52\% error reduction) and Brier score (0.013 vs.\ 0.024, a 46\% improvement). A model size ablation reveals that the joint model's generation loss acts as an implicit calibration regularizer, while specialization primarily improves the sharpness of predictions. These results suggest that practitioners seeking reliable copy decisions should train a specialized model---not because it improves calibration, but because it produces substantially more accurate and confident predictions.